\documentclass[12pt]{article}
\usepackage{amsmath,amssymb,booktabs,longtable,array}
\usepackage{fontspec}
\usepackage{polyglossia}
\setmainlanguage{arabic}
\setotherlanguage{english}
\newfontfamily\arabicfont[Script=Arabic]{Amiri}
\usepackage{geometry}
\geometry{margin=2.1cm}
\usepackage{hyperref}

\title{اختبار Scaling عبر أكثر من خمس مراتب في \(X\)}
\author{مشروع تجربة البحث عن الأعداد الأولية}
\date{2026-09-01}

\begin{document}
\maketitle

\section{السؤال}

في التجربة السابقة ظهر أن التمثيل المحلي متعدد المقاييس يصبح أكثر ثباتًا
عندما نستعمل المتغير البعدي:
\[
L/H
\]
بدل طول الدعم المطلق \(L\).

لكن نطاقنا كان صغيرًا نسبيًا:
\[
X\approx1.6\times10^7
\text{ إلى }
6.3\times10^7.
\]

السؤال الحالي هو:

\begin{quote}
هل يظل الطيف البعدي نفسه تقريبًا عندما نرفع \(X\) عبر عدة مراتب عشرية،
أم أن collapse السابق كان أثر نطاق محدود؟
\end{quote}

\section{تنبيه مهم: ما الذي نختبره هنا؟}

هذه التجربة لا تعد الأوليات الفعلية حتى \(10^{12}\).

نحن نولد هندسة فترات:
\[
p_n^2<x<p_{n+1}^2
\]
عند قيم أكبر من \(p_n\)، ثم نبني target covariance من قانون الفترات القصيرة
نفسه الذي استعملناه سابقًا.

إذن هذه تجربة على:

\[
\boxed{
\text{استقرار تمثيلنا المحلي للنظرية على هندسة prime-square حقيقية عند }X\text{ كبير.}
}
\]

وليست بعد اختبارًا تجريبيًا مباشرًا لأعداد الأوليات الفعلية عند \(10^{12}\).

\section{النطاق}

اخترنا ستة مستويات تقريبية:

\[
2\times10^7,\quad
2\times10^8,\quad
2\times10^9,\quad
2\times10^{10},\quad
2\times10^{11},\quad
2\times10^{12}.
\]

حول كل مستوى أخذنا ثلاث لوحات من \(128\) فترة متتالية بين مربعات أوليات،
أي:

\[
18
\]
لوحة إجمالًا.

النطاق الحقيقي للـmedian \(X\) كان:

\[
1.17\times10^7
\le X\le
2.01\times10^{12},
\]
أي نحو:

\[
\boxed{5.24\text{ مراتب عشرية}.}
\]

\section{تحسين منهجي في الدقة}

في التجربة السابقة كان حجم الخلية ثابتًا:
\[
1024.
\]

هذا غير مناسب عند تغيير \(H\) بمئات المرات.

لذلك جعلنا دقة الشبكة نفسها scale-covariant:

\[
\mathrm{cell}
\approx
\frac{H_{\rm median}}{64}.
\]

أي أننا نحافظ تقريبًا على \(64\) خلية في طول فترة نموذجي.

كما استعملنا شبكة مقاييس كثيفة:
\[
16
\]
مقياسًا في كل octave.

وبذلك لا يمكن أن يكون ثبات \(L/H\) مجرد أثر من دقة مطلقة ثابتة.

\section{لماذا جمعنا النطاقات أكثر؟}

أعدنا تأكيد مشكلة identifiability:

المعاملات الفردية بين المقاييس القريبة يمكن أن تتبادل الوزن لأن covariance
التي تولدها تكون شديدة التشابه.

لذلك استعملنا هذه النطاقات الخشنة:

\[
L/H<4,
\]
\[
4\le L/H<8,
\]
\[
L/H\ge8.
\]

هذه أكثر ثباتًا من coefficient منفرد أو حتى octave ضيق.

\section{النتيجة الأساسية}

عبر جميع اللوحات الثماني عشرة، كانت مساهمة:

\[
L/H<4
\]
هي:

\[
\boxed{
0.8426\pm0.0410.
}
\]

وكان مجالها الكامل:
\[
0.7540
\text{ إلى }
0.8935.
\]

أما:

\[
4\le L/H<8
\]
فكان:

\[
0.0914\pm0.0411.
\]

وأما:

\[
L/H\ge8
\]
فكان:

\[
\boxed{
0.0660\pm0.0132.
}
\]

إذن، رغم زيادة \(X\) بأكثر من \(10^5\)، بقي الطيف الخشن قريبًا جدًا من
نفس التقسيم:

\[
\boxed{
\text{نحو }84\%\text{ تحت }4H,
\quad
9\%\text{ بين }4H\text{ و }8H,
\quad
7\%\text{ فوق }8H.
}
\]

\section{هل يوجد drift مع } \(X\)؟

أجرينا انحدارًا خطيًا على:
\[
\log_{10}X.
\]

للنطاق:
\[
L/H<4
\]
كانت الميل لكل decade:

\[
0.00303\pm0.00577,
\]
مع:
\[
p=0.606.
\]

للنطاق:
\[
4\le L/H<8
\]
كان:

\[
-0.00576\pm0.00564,
\]
مع:
\[
p=0.322.
\]

وللنطاق:
\[
L/H\ge8
\]
كان:

\[
0.00273\pm0.00175,
\]
مع:
\[
p=0.138.
\]

إذن لا يوجد في هذه العينة دليل إحصائي على drift خطي واضح عبر خمسة decades.

وهذا لا يثبت invariance مطلقة؛ لكنه يرفض فكرة أن collapse السابق كان
ينهار سريعًا عندما نكبر \(X\).

\section{السلوك عند كل مستوى}

متوسط مساهمة:
\[
L/H<4
\]
عند مستويات \(X\) الستة كان تقريبًا:

\[
0.831,\quad
0.831,\quad
0.835,\quad
0.876,\quad
0.847,\quad
0.835.
\]

أما:
\[
L/H\ge8
\]
فكان:

\[
0.0557,\quad
0.0687,\quad
0.0710,\quad
0.0536,\quad
0.0683,\quad
0.0788.
\]

إذن التقلب بين اللوحات أكبر من أي اتجاه أحادي واضح مع \(X\).

\section{اختبار amplitude النظري}

في الحالة المثالية ذات الفترات المتساوية، تعطي نظرية Leung تقريبًا:

\[
\rho_1
\sim
-\frac{\log2}{\log(X/H)}.
\]

لذلك فحصنا:
\[
\rho_1\log(X/H).
\]

عبر اللوحات الثماني عشرة حصلنا على:

\[
\boxed{
-0.6523\pm0.0111.
}
\]

بينما:
\[
-\log2=-0.6931.
\]

القيمة ليست مساوية تمامًا لـ\(-\log2\)، لأن فترات مربعات الأوليات غير
متساوية ولأننا نستعمل صيغة variable-interval covariance لا الحالة المثالية.

لكن التشتت النسبي حول المتوسط صغير، نحو بضعة بالمئة فقط، عبر أكثر من خمس
مراتب في \(X\).

إذن فصل:

\[
\boxed{
\text{shape dimensionless}
+
\text{logarithmic amplitude}
}
\]
يبقى وصفًا جيدًا على هذا النطاق.

\section{محاولة أقوى فشلت: spectrum دقيق واحد للجميع}

حاولنا أخذ fine-grained spectrum واحد من حالة مرجعية متساوية الفترات،
ثم نقله حرفيًا إلى prime-square panels باستعمال \(L/H\) وfactor لوغاريتمي
للـamplitude.

هذه المحاولة فشلت.

متوسط kernel RMSE كان:

\[
0.02595,
\]
وأفضل حالة فقط:

\[
0.01444.
\]

وفي بعض اللوحات أنتج النقل الدقيق حتى إشارة خاطئة عند lag 1.

إذن:

\[
\boxed{
\text{الطيف الدقيق ليس universal إذا استعملنا median }H\text{ واحدًا للوحة.}
}
\]

هذه نتيجة سلبية مهمة.

السبب المرجح هو أن prime-square intervals داخل اللوحة نفسها تملك:
\[
H_i
\]
شديدة التفاوت بسبب prime gaps.

أي أن استعمال:
\[
L/H_{\rm median}
\]
جيد للطيف الخشن، لكنه غير كافٍ لنقل كل coefficient الدقيق.

\section{ما الذي ثبت وما الذي لم يثبت؟}

الذي تدعمه التجربة بقوة هو:

\[
\boxed{
\text{coarse dimensionless spectrum مستقر تقريبًا عبر }5.2\text{ decades في }X.
}
\]

أما العبارة الأقوى:

\[
\text{يوجد fine spectrum واحد }W(L/H)\text{ يصلح لكل prime-square panel}
\]

فلم تدعمها التجربة.

إذن الصياغة الأدق أصبحت:

\[
\boxed{
\mathcal W_{\rm coarse}(L/H)
\text{ مستقرة نسبيًا،}
}
\]
بينما البنية الدقيقة تحتاج إلى معرفة:
\[
H_i
\]
المحلي، لا median \(H\) فقط.

\section{تفسير ما تعلمناه}

هذه النتيجة تعيد توجيه الآلية المحلية.

بدل أن نضع موجات Haar على خط \(n\) بمقياس عالمي داخل اللوحة، قد نحتاج
إلى coordinate محلي متكيف، بحيث تطبع كل فترة بطولها الخاص:

\[
u
=
\int \frac{dn}{H_{\rm local}(n)}.
\]

في أبسط صورة، نجعل كل prime-square interval وحدة طول واحدة في coordinate
الجديد.

ثم نسأل:

\begin{quote}
هل يظهر fine spectrum universal عندما نقيس الدعم في هذا coordinate
المتكيف بدل استعمال طول فعلي واحد \(L\) مقسومًا على median \(H\)؟
\end{quote}

هذا هو الاختبار الطبيعي التالي الذي ولدته نتيجة الفشل نفسها.

\section{الحكم}

نتيجة الـscaling ليست ``قانونًا جديدًا للأوليات''، لأن target kernel نفسه
مأخوذ من نظرية الفترات القصيرة.

لكنها نتيجة مهمة عن نموذجنا:

\[
\boxed{
\text{الآلية المحلية لا تنهار عندما ننقلها من }10^7\text{ إلى }10^{12}.
}
\]

وما يبقى ثابتًا ليس coefficient بعينه، بل توزيع خشن لمساهمة covariance
بحسب المقياس النسبي \(L/H\).

في المقابل، فشل النقل الدقيق كشف لنا المتغير المفقود التالي:

\[
\boxed{H_i\text{ المحلي داخل اللوحة}.}
\]

\section{ملفات التجربة}

\begin{itemize}
\item \texttt{wide\_X\_multiscale\_scaling\_panels.csv}
\item \texttt{wide\_X\_multiscale\_scaling\_grouped.csv}
\item \texttt{wide\_X\_multiscale\_scaling\_regression.csv}
\item \texttt{wide\_X\_rho1\_log\_scaling.csv}
\item \texttt{wide\_X\_multiscale\_scaling\_summary.csv}
\item \texttt{wide\_X\_universal\_spectrum\_transfer.csv}
\item \texttt{analyze\_wide\_X\_multiscale\_scaling.py}
\item \texttt{plots/wide\_X\_multiscale\_coarse\_bands.svg}
\item \texttt{plots/wide\_X\_rho1\_log\_scaling.svg}
\item \texttt{plots/wide\_X\_fine\_spectrum\_transfer\_failure.svg}
\end{itemize}

\end{document}
